\chapter{Implementation Code}
\label{app:code}

This appendix describes the implementation and how the results in this dissertation can be reproduced. The full source is available in the repository \texttt{\seqsplit{ryantjx/sqmc-cuthbert}} (currently a draft) and will be integrated into \texttt{state-space-models/cuthbert} under the issue \href{https://github.com/state-space-models/cuthbert/issues/255}{Implement Sequential Quasi Monte Carlo \#255}. Figure~\ref{fig:issue-255} shows the tracking issue. The repository is organised as two packages: \texttt{sqmc/}, containing the shared SQMC library (QMC point generation, the SQMC filter and Hilbert sorting) together with its CPU/GPU benchmark suite, and \texttt{rbsqmc/}, containing the football application (data preparation, the RB-SMC, RB-SQMC and factorial EKF filters, and the comparison pipelines). The football package imports the QMC generation and Hilbert-sorting modules from the SQMC package at runtime. A replication guide with the complete commands is provided in the repository's \texttt{README.md}; this section summarises the steps.

\section{Replication}
\label{app:replication}

The experiments require Python~3.10 or later with the dependencies listed in the repository's \texttt{requirements.txt}. GPU runs use the CUDA~12 build of \texttt{jax}; the football code selects its platform through the \texttt{RBSQMC\_PLATFORM} environment variable. Two data files are generated on a fresh clone: the Sobol' direction-number table, built from the committed Joe--Kuo source table by \texttt{\_generate\_sobol\_data.py}, and the parquet cache of the international football results, downloaded from the \texttt{martj42/international-results} source. All commands are run from the repository root.

Table~\ref{tab:replication} maps each empirical chapter to the entry point that reproduces it. The performance comparison of Chapter~\ref{ch:sqmc} is produced by three standalone benchmarks, which can be run directly on a CPU or GPU host or through the Colab launcher, which provisions an A100 session, executes the three stages sequentially and verifies each stage's artefacts before proceeding. The football comparison of Chapter~\ref{ch:model} trains the correlated RB-SQMC model and the factorial EKF baseline on identical data splits and optimiser settings under the configuration recorded in Appendix~\ref{app:config}, and evaluates one-step-ahead predictions on the held-out period. Each run writes a timestamped output directory containing the effective configuration, device and software metadata, the git commit, source hashes, results and figures, so that any reported number can be traced to its generating configuration.

\begin{table}[H]
\centering
\begin{tabularx}{\textwidth}{lXX}
\toprule
Result & Entry point & Output \\
\midrule
QMC generation (Sec.~\ref{sec:qmc-benchmark}) & \texttt{\seqsplit{sqmc.comparison.benchmark\_qmc}} & \texttt{\seqsplit{sqmc/comparison/outputs/}} \\
Hilbert sorting (Sec.~\ref{sec:hilbert-benchmark}) & \texttt{\seqsplit{sqmc.comparison.benchmark\_hilbert\_sort}} & \texttt{\seqsplit{sqmc/comparison/outputs/}} \\
SQMC CPU/GPU (Sec.~\ref{sec:sqmc-cpu-gpu}) & \texttt{\seqsplit{sqmc.comparison.benchmark\_sqmc}} & \texttt{\seqsplit{sqmc/comparison/outputs/}} \\
RB-SQMC vs EKF (Sec.~\ref{sec:performance-comparison}) & \texttt{\seqsplit{rbsqmc.comparison.sqmc\_ekf.run}} & \texttt{\seqsplit{rbsqmc/comparison/sqmc\_ekf/outputs/}} \\
RB-SMC vs RB-SQMC & \texttt{\seqsplit{rbsqmc.comparison.sqmc\_smc.compare\_smc\_sqmc}} & \texttt{\seqsplit{rbsqmc/outputs/compare/}} \\
\bottomrule
\end{tabularx}
\caption{Entry points reproducing the empirical results. Each writes a timestamped directory with configuration, provenance metadata and results.}
\label{tab:replication}
\end{table}

The authoritative football run is \texttt{09092026\_1431} at source commit \texttt{336842f} on an NVIDIA A100, as recorded in Appendix~\ref{app:config}. The benchmark suite records the same provenance metadata (device descriptions, software versions, git commit and hashes of the benchmark sources and the Sobol' direction data) in every run directory. Timed measurements use seven repetitions after two warm-up calls, with compilation and input transfers excluded from the reported runtimes; the SQMC accuracy comparison selects particle counts under a per-trajectory runtime budget using selection randomisations and reports held-out accuracy from separate validation randomisations, as described in Section~\ref{sec:empirical-design}.

\begin{figure}[t]
\centering
\includegraphics[width=0.9\textwidth]{https-github.com-state-space-models-cuthbert-issues-255.png}
\caption{The tracking issue for the SQMC implementation in the \texttt{state-space-models/cuthbert} repository.}
\label{fig:issue-255}
\end{figure}

\section{Quasi-Monte Carlo point generation (\texttt{qmc.py})}
\label{app:qmc-code}

The following is the core of the QMC module, which implements the Halton and Sobol' sequences described in Section~\ref{sec:qmc-sequences}. The full file is in \texttt{sqmc/qmc/qmc.py}.

% \lstinputlisting[language=Python, caption={The QMC module implementing the Halton and Sobol' sequences.}, label={lst:qmc}]{../../sqmc/qmc/qmc.py}

\section{Hilbert sorting (\texttt{hilbert\_sort.py})}
\label{app:hilbert-code}

The following is the Hilbert-sort module, which computes the packed Hilbert index of each particle and sorts them, as described in Section~\ref{sec:sqmc}. The full file is in \texttt{sqmc/hilbert\_sort/hilbert\_sort.py}.

% \lstinputlisting[language=Python, caption={The Hilbert-sort module computing packed Hilbert indices and sorting particles.}, label={lst:hilbert}]{../../sqmc/hilbert_sort/hilbert_sort.py}
